
\documentclass[11pt]{article}
\usepackage[a4paper,margin=1in]{geometry}
\usepackage{amsmath,amssymb,mathtools}
\usepackage{booktabs}
\usepackage{longtable}
\usepackage{array}
\usepackage{hyperref}
\usepackage{enumitem}
\usepackage[T1]{fontenc}
\usepackage[utf8]{inputenc}

\title{Origins/NOEMA Semantic--Nucleic Braiding Protocol v0.8\\
\large Internet-backed Comparisons, Citations, and Model Positioning}
\author{CIEL/NOEMA working artifact for Adrian Lipa}
\date{2026-05-17}

\begin{document}
\maketitle

\begin{abstract}
This v0.8 document preserves the deterministic v0.5--v0.7 NOEMA model while adding internet-backed comparisons to established origin-of-life theories. The model remains a proxy model: it does not claim empirical proof of abiogenesis, quantum chemistry, a thermodynamic membrane simulation, or a geological timeline. Its contribution is narrower: it formalises how a four-block nucleic grammar can be coupled to time-varying environmental forcings and tested for invariant retention, transition readiness, and protocell-like threshold crossing. The comparison layer situates this proxy against RNA world, warm-little-pond, mineral templating, iron--sulfur/metabolism-first, alkaline hydrothermal vent, lipid/protocell, condensate, and classical primordial-soup models.
\end{abstract}

\section*{NOEMA status}
\begin{verbatim}
sot_available: local_minimal_noema_from_uploaded_artifacts
current_state_loaded: true
current_state_id: NOEMA-20260517T061500-v08-internet-comparisons-citations
previous_state_id: NOEMA-20260517T052000-v07-temporal-climate-options
memory_layer: episodic
uncertainty_operator: chyba
canon_allowed: false
pdf_generated: false
internet_used_for_references: true
\end{verbatim}

\section{Scope and epistemic boundary}
The v0.8 model compares the internal NOEMA proxy with established origin-of-life theories rather than replacing them. The model's ``DNA four-block'' grammar is not a claim that modern DNA preceded RNA. It is a relational alphabet proxy inherited from the Watson--Crick complementarity pattern \cite{WatsonCrick1953} and used to measure complementarity, symmetry, compression, and protocell-readiness. This must be distinguished from the RNA world hypothesis, where RNA is considered a plausible earlier carrier of both genetic information and catalytic activity \cite{RobertsonJoyce2012}.

Classical origin-of-life research includes the Oparin--Haldane primordial-soup tradition \cite{Oparin1938,Haldane1929}, the Miller--Urey demonstration of abiotic amino-acid synthesis under a then-plausible reducing atmosphere \cite{Miller1953}, metabolism-first and iron--sulfur surface-metabolism models \cite{Wachtershauser1988}, RNA-world and prebiotic nucleotide-synthesis programmes \cite{PownerGerlandSutherland2009,RobertsonJoyce2012}, hot-spring/wet--dry-cycle scenarios \cite{DamerDeamer2015,DamerDeamer2020}, alkaline-vent models \cite{MartinRussell2003,MartinBarossKelleyRussell2008}, and protocell membrane models \cite{Mansy2008Protocells}. The NOEMA model is positioned as a computational bridge between these domains: it represents environmental variation as time-dependent forcing and evaluates whether a relational information core retains enough invariant structure to cross a protocell-readiness threshold.

\section{Existing project metrics retained from v0.5--v0.7}
The v0.8 comparison layer does not alter the v0.5/v0.6/v0.7 numeric core. It retains the four-block DNA relation metrics and the v0.7 temporal option results.

\begin{longtable}{@{}ll@{}}
\toprule
Metric & Value \\
\midrule
Threshold proxy & 0.692709 \\
Baseline peak readiness & 0.693658 \\
Best option & ionic_moderation_window \\
Best option peak readiness & 0.708705 \\
Fastest growth option & wet_dry_pulse_amplified \\
Fastest growth maximum velocity & 0.256095 \\
Negative control & uv_and_dilution_stress \\
Negative control peak readiness & 0.670506 \\
Negative control transition detected & false \\
\bottomrule
\end{longtable}

\section{Comparison 1: four-block grammar versus RNA-world models}
The NOEMA four-block grammar compresses A, C, G, and T into complementarity axes, purine--pyrimidine class symmetry, and a stability gradient. Classical RNA-world models instead ask whether RNA-like molecules could both store information and catalyse reactions before DNA/protein systems emerged \cite{RobertsonJoyce2012}. Sutherland's activated pyrimidine ribonucleotide synthesis showed that prebiotic pathways need not assemble sugar, base, and phosphate by naive stepwise addition; they can proceed through chemically plausible precursor networks \cite{PownerGerlandSutherland2009}.

The comparison is therefore not one of contradiction. The NOEMA grammar is a higher-level relation model; RNA-world chemistry is a molecular feasibility model. The grammar tells us whether complementarity and class closure form a compressible information core. RNA-world chemistry tells us whether the relevant molecular substrates can plausibly arise and participate in catalysis or templating.

\section{Comparison 2: semantic compression versus hypercycles and quasispecies}
Eigen's self-organisation and hypercycle framework treats early molecular evolution as the emergence of self-replicating information-bearing macromolecules under mutation, selection, and catalytic coupling \cite{Eigen1971,EigenSchuster1979}. The NOEMA semantic-compression layer is not a population-genetic model. It does not model mutation spectra, replication fidelity, or error thresholds directly. Instead, it measures whether the relational trace of a process remains identifiable after surface reduction.

This makes the NOEMA metric complementary to hypercycle/quasispecies approaches. Hypercycle models ask whether replicator networks can persist under error and selection. NOEMA asks whether the process has an invariant relational core at all: recurrence, residue memory, transition graph continuity, and threshold-crossing behaviour. If future work adds explicit mutation and replication operators, the v0.8 compression trace can become the state representation on which quasispecies-style dynamics are run.

\section{Comparison 3: hydration cycles and warm-little-pond/hot-spring models}
The v0.7 forcing option \texttt{wet\_dry\_pulse\_amplified} produced the fastest growth velocity. This aligns structurally with warm-little-pond and hot-spring scenarios, where hydration/dehydration cycles concentrate solutes, promote condensation reactions, and allow repeated selection over partially assembled products \cite{Pearce2017WarmPonds,DamerDeamer2015,DamerDeamer2020}. Powner, Gerland, and Sutherland's nucleotide synthesis also used sequential aqueous, heating, evaporation, irradiation, and phosphorylation steps, making cyclic environmental forcing chemically relevant rather than merely metaphorical \cite{PownerGerlandSutherland2009}.

The NOEMA model converts this into an explicit option class: a hydration cycle is represented as a time-dependent pulse on concentration, recurrence, and compartmentalisation pressure. The model therefore agrees with wet--dry scenarios at the level of process topology: cyclic forcing can increase readiness more effectively than static exposure.

\section{Comparison 4: mineral templating, clay catalysis, and surface metabolism}
Mineral templating in NOEMA is treated as a structured surface forcing field. This can be compared with montmorillonite-catalysed RNA oligomer formation and mineral-surface prebiotic oligomer synthesis \cite{FerrisErtem1993Montmorillonite,Ferris1996MineralSurfaces}. It also overlaps conceptually with iron--sulfur surface-metabolism models, where mineral surfaces are not passive containers but catalytic and organisational participants \cite{Wachtershauser1988}.

The NOEMA distinction is that mineral templating is not assumed to be sufficient. In the v0.7/v0.8 framework it is one component among hydration cycling, thermal gradients, ionic moderation, amphiphile availability, compartmentalisation pressure, and holonomic recurrence. This prevents a single-surface explanation from becoming an untested monofactor model.

\section{Comparison 5: thermal gradients and alkaline hydrothermal vent theories}
Alkaline hydrothermal vent models emphasise sustained geochemical disequilibrium, microporous compartments, pH gradients, temperature gradients, and mineral membranes as plausible drivers of early metabolism and cellularisation \cite{MartinRussell2003,MartinBarossKelleyRussell2008}. The v0.7 thermal-gradient field is compatible with that logic, but it remains environment-agnostic: it can represent a vent gradient, a geothermal pool gradient, or a local surface gradient without committing to one geological site.

The NOEMA model therefore treats thermal gradients as a rate-and-selection operator. It asks whether a gradient contributes to increasing readiness, stabilising transition traces, or producing destructive stress. This is why the negative control \texttt{uv\_and\_dilution\_stress} fails to cross the threshold: energy input alone is not enough if it does not preserve recurrence, compartmentalisation, and residue continuity.

\section{Comparison 6: ionic moderation, amphiphiles, and lipid/protocell models}
The best v0.7 option was \texttt{ionic\_moderation\_window}, with a peak readiness proxy of 0.708705. This is consistent with protocell membrane work where amphiphile assembly, vesicle stability, permeability, and ionic context jointly constrain primitive compartment viability \cite{Mansy2008Protocells}. Fatty-acid protocell systems are dynamically different from modern phospholipid membranes, and recent model-protocell work also treats lipid exchange and fusion as important dynamical constraints \cite{Fan2023LipidExchange}.

In NOEMA terms, ionic moderation is not merely a chemical parameter; it is a boundary-condition stabiliser. If the medium is too dilute, the system cannot retain local reaction products. If too harsh or salt-loaded, compartment formation and polymer compatibility are compromised. The option window formalism is therefore appropriate: it models a habitable band rather than a one-directional increase.

\section{Comparison 7: compartmentalisation pressure versus vesicles, droplets, condensates, and vents}
Compartmentalisation is the strongest recurring driver in the v0.6/v0.7 stage tables. This aligns with multiple theories: fatty-acid vesicles as protocell containers \cite{Mansy2008Protocells}, vent micropores as geochemical compartments \cite{MartinRussell2003}, hydrothermal pool vesicle encapsulation \cite{DamerDeamer2015}, and more recent condensate-centric RNA models \cite{FineMoses2024Condensate}.

The NOEMA model abstracts these into a common operator: co-localisation plus boundary memory. A compartment is useful only if it changes the retention of reaction products and histories. In this sense the model unifies membrane-first, metabolism-first compartment, droplet/condensate, and vent-pore scenarios at the level of relational topology.

\section{Comparison 8: holonomic recurrence versus autocatalytic and systems-chemistry views}
Holonomic recurrence is the NOEMA field that measures whether the system re-enters similar relational states while preserving enough memory of prior transitions. It is comparable to autocatalytic network persistence, systems chemistry, and selection over coupled reaction networks \cite{Szostak2009,EigenSchuster1979}. However, NOEMA does not yet implement full reaction kinetics. It implements a recurrence-and-retention proxy.

This means holonomic recurrence should be interpreted as a diagnostic layer: it detects whether environmental cycling and compartmentalisation produce a memory-bearing trace. It does not yet prove autocatalysis, heredity, or Darwinian evolution. Those require future explicit reaction-network and replication modules.

\section{Comparison 9: climate-holonomic change through time}
The v0.7 model treats climate-holonomic change as relative-time forcing, not as an empirical Hadean chronology. This is consistent with calls to embed prebiotic chemistry in changing environmental and planetary contexts rather than treating it as a fixed laboratory recipe \cite{Lyons2020Environments}. Warm ponds, vents, hot springs, impacts, UV exposure, minerals, salinity, and atmospheric conditions each change the feasible reaction space.

The current model therefore supports scenario testing. A user can encode environmental options as offsets, slopes, and pulses:
\begin{verbatim}
TemporalForcingOption(
  field_offsets={...},
  field_slopes={...},
  field_pulses={...}
)
\end{verbatim}
This is the correct computational interface for the next stage: tempo and kind of change can be varied while the v0.5 information-core metrics remain fixed.

\section{Results and interpretation}
The model's strongest result remains that the transition is front-loaded. The largest readiness jump occurs when an information scaffold is first coupled to concentration, surface cycling, and compartmentalisation. Subsequent growth is slower and stabilises near the threshold. The best option is an ionic moderation window, while wet--dry pulsing gives the fastest growth rate. The negative control fails because destructive stress and dilution do not preserve enough recurrence and compartmental trace.

The result is compatible with several existing theories but identical to none. It is not RNA-world chemistry, although it can host RNA-world variables. It is not a vent model, although it can encode vent-like gradients. It is not a lipid-world model, although it treats amphiphile availability and compartment pressure explicitly. Its distinctive contribution is the invariant-trace criterion: a protocell candidate appears when information grammar, environmental cycling, and boundary formation preserve a process identity over time.

\section{Limitations and next technical steps}
The v0.8 model remains a deterministic proxy. It does not simulate molecular dynamics, quantum chemistry, membrane thermodynamics, explicit reaction kinetics, mutation, replication fidelity, or real geochemical chronology. The correct next steps are:
\begin{enumerate}[leftmargin=*]
  \item Add explicit reaction-network modules for precursor formation, oligomerisation, hydrolysis, and autocatalytic closure.
  \item Add environmental parameter ranges tied to literature-backed plausible Hadean settings.
  \item Add sensitivity analysis over hydration period, salt concentration, amphiphile abundance, mineral-surface affinity, UV stress, dilution, and thermal amplitude.
  \item Add a separate validation table distinguishing empirical support, model analogy, and pure NOEMA proxy assumptions.
  \item Preserve the NOEMA append-only workflow: every change requires a separate update snap and a separate results package.
\end{enumerate}

\section{Epistemic status}
\begin{verbatim}
semantic_relation_model_claimed: true
time_varying_environment_options_claimed: true
internet_backed_comparisons_added: true
bibtex_references_added: true
empirical_abiogenesis_proof_claimed: false
DNA_first_claimed: false
quantum_chemistry_claimed: false
thermodynamic_membrane_simulation_claimed: false
geological_timeline_claimed: false
pdf_generated: false
\end{verbatim}

\bibliographystyle{plain}
\bibliography{references}

\end{document}
